\begin{scanpage}{0112}{95}
disk $(\omega=\frac\pi2,\ \Theta=\pi)$, the field from the nonuniform part of the current equals zero on the $z$ axis, since $f^1=-g^1=-1/2$ when $\vartheta=0$, and $f^1=g^1=-1/2$ when $\vartheta=\pi$ [compare (8.16)].

Using the resulting relationships in the following sections, we will calculate the effective scattering area (in the direction $\vartheta=\pi$) for specific bodies. We shall assume that they are irradiated by the plane wave
\[E_x=H_y=E_{0x}e^{ikz},\tag{16.20}\]
and their linear dimensions are large in comparison with the wavelength.
\sectiontitle{17}{A Cone}
Let a cone (Figure 28) be irradiated by plane electromagnetic wave (16.20). The uniform part of the current which is excited on the cone's surface has the components
\[\left.\begin{aligned}j_x^0&=\frac c{2\pi}E_{0x}\sin\omega\,e^{ikz},\\j_y^0&=0,\\j_z^0&=\frac c{2\pi}E_{0x}\cos\omega\cos\psi\,e^{ikz}\end{aligned}\right\}\tag{17.01}\]
and creates in the direction $\vartheta=\pi$ (with $R\gg ka^2$, $R\gg kl^2$) the field
\[\left.\begin{aligned}E_x=-H_y&=-E_{0x}\frac i{4k}\operatorname{tg}^2\omega\frac{e^{ikR}}R\\&\quad+E_{0x}\left(\frac i{4k}\operatorname{tg}^2\omega+\frac a2\operatorname{tg}\omega\right)\frac{e^{ikR}}R e^{2ikl},\\E_y=H_x&=0.
\end{aligned}\right\}\tag{17.02}\]
Here the first term describes the spherical wave diverging from the vertex of the cone, and the remaining terms describe the spherical wave from its base.

The field caused by the discontinuity of the surface at the cone base is a spherical wave, and is determined in accordance with
\end{scanpage}
